\begin{table}[H]
\centering
\small
\setlength{\tabcolsep}{4pt}
\begin{tabular}{lc|cc|cc}
\toprule
& & \multicolumn{2}{c|}{spectrum null A} & \multicolumn{2}{c}{hub-randomizing null B} \\
synthetic cloud ($n{=}1000$, $d{=}768$) & $\hat\delta$ & Gaussian & Haar & Gaussian & Haar \\
\midrule
30-cluster star, tight (0.1) & $0.055$ & $-0.115$ & $-0.115$ & $-0.126$ & $-0.058$ \\
30-cluster star, mid (0.3) & $0.059$ & $-0.104$ & $-0.103$ & $-0.114$ & $-0.049$ \\
30-cluster star, loose (0.6) & $0.061$ & $-0.084$ & $-0.082$ & $-0.090$ & $-0.033$ \\
2-level hierarchy $6\times5$ (0.3) & $0.044$ & $-0.143$ & $-0.143$ & $-0.113$ & $-0.112$ \\
\bottomrule
\end{tabular}
\caption{Calibration of the nulls on synthetic clouds (excess, means over seeds; within/between-cluster noise ratio in parentheses). A pure star (Gaussian clusters around Gaussian centers, no hierarchy) sits far below the spectrum null A under either construction (Gaussian coefficients, the paper's; Haar-rotated coefficients, exact sample spectrum): null A certifies clustering, not depth. The hub-randomizing null B, which keeps every cluster intact and resamples only the hub configuration, also reports the star as ``hierarchical'': a near-regular simplex of hubs already minimizes $\delta$, so any hub randomization raises it. Null B is therefore not a valid depth test as constructed and is not used in the paper; a star-calibrated version (excess B relative to a matched star with the same cluster count and tightness) may still serve as a depth test, since under the Haar construction the two-level hierarchy sits $2$--$3\times$ below the star. That version is run on the real backbones in Table~\ref{tab:b34-depthvariants}. Under the plain Gaussian null the mid-radius star scores $-0.104$ and the $6{\times}5$ hierarchy $-0.14$: a star and a two-level tree differ by a fraction of either's excess, which is why depth cannot be read from the spectrum excess alone.}
\label{tab:b25-star}
\end{table}
